\section{How many Jensen tests can a sextic perturbation pass?}
\label{sec:kernel}
The incoming report \emph{Geometric progressions and nonabelian zero
tests} treats the conductor-$23$ Artin factor through the Mellin kernel of
$\eta(z)\eta(23z)$. Its sections on Jensen polynomials give all shifts
through degree $16$ by an independent finite-height certificate, and
through degree $945{,}173$ using \href{https://www.cs.uleth.ca/~nathanng/RESEARCH/phd.thesis.pdf}{Ng's published zero count}. Its separate
counterexample for $\exp(-u^2-\delta u^6)$ shows why extending those
conclusions to every degree from kernel curvature alone is problematic.
We investigate a question suggested by that counterexample. Can the
quartic term restore every low-degree test while a higher-degree test
still detects nonreal zeros?

Fix $t>0$ and $c\ge0$, and consider
\[
 K_{t,c}(u)=\exp(-u^2-t^2u^4-ct^3u^6).
\]
Write $\int_{\mathbb R}e^{zu}K_{t,c}(u)\,du
=\sum_{n\ge0}\gamma_n z^{2n}/n!$, and put
$b_n=4^n\gamma_n/\sqrt\pi$.
Scaling the coefficients this way simply rescales the Jensen variable
and multiplies the polynomial by a positive constant.
The degree-$d$, shift-$n$ polynomial is
$J_{d,n}(Z)=\sum_{k=0}^d\binom dk b_{n+k}Z^k$.
Here real-rooted means that all roots, counted with multiplicity, are real.

\paragraph{The scaling limit.}
For each fixed $n,d,c$, with $D=d/dy$,
\[
 t^{-d}J_{d,n}(-1+ty)\ \longrightarrow\ Q_{d,c}(y)
 =\exp(-D^2+cD^3)y^d
 \qquad(t\downarrow0)
\]
coefficientwise. The exponential on the right is a finite sum on
polynomials of bounded degree. The model polynomials themselves are
classical three-variable Hermite polynomials \cite{hermite}. The limit identifies
which ones control this two-parameter family of Mellin moments.

\emph{Proof.}
Writing $p=n+1/2$ and $X\sim\mathrm{Gamma}(p,1)$ gives
$b_n=\mathbb E H(tX)$, where $H(s)=\exp(-s^2-cs^3)$, by $x=u^2$
and the duplication formula. If $h_j=[s^j]H(s)$, Taylor's formula gives
\[
 b_n=\sum_{j=0}^d h_j t^j (p)_j+O_{n,d,c}(t^{d+1}).
\]
Here $(p)_j$ is the rising factorial. Every derivative of $H$ is bounded
on the positive half-line and $X$ has every moment. At index $n+k$,
$(p+k)_j$ is monic of degree $j$ in $k$. Acting on $(1+Z)^d$ with
$(Z\partial_Z)^j$ inserts the multiplier $k^j$ in its binomial sum.
The top derivative in that operator is $Z^j\partial_Z^j$. At $Z=-1+ty$
this yields $(-1)^j d^{\underline j}t^{d-j}y^{d-j}+O(t^{d-j+1})$,
where $d^{\underline j}=d!/(d-j)!$.
All lower derivatives have one extra power of $t$ or more.
Summing gives $H(-D)y^d$. The remainders are coefficientwise since
the degree is bounded and only finitely many moments are involved.

\paragraph{The sharp cubic transition.}
Put $c_*=2\sqrt2/3$. For every fixed shift the cubic has three distinct
real roots for all sufficiently small $t>0$ if $0\le c\le c_*$,
and has a nonreal pair for all sufficiently small $t>0$ if $c>c_*$.
Indeed, $Q_{3,c}=y^3-6y+6c$ has discriminant $864-972c^2$. Away
from the boundary the assertion follows by continuity of the roots.
At $c_*$ the limiting polynomial has a double root at $y=\sqrt2$.
Using the first five Taylor coefficients $(1,0,-1,-c,1/2)$ in the
finite binomial sum gives
\[
 t^{-3}J_{3,n}(-1+ty)=Q_{3,c}(y)+6t(p+3)(y^2-3cy-2)+O_{n,c}(t^2).
\]
The correction at the double root is $-24t(p+3)$.
The local minimum is therefore negative and the values a fixed small
distance to either side are positive. The simple third root also persists.
More precisely, at $c=c_*+at$ the boundary value $a=4(p+3)$ separates
the two strict outcomes to this order: the local minimum is
$t(6a-24(p+3))+O(t^2)$, since the critical point moves by $O(t)$.
All these statements keep the shift fixed.

\paragraph{Every positive sextic coefficient is visible at finite degree.}
If $d\ge3$ and $Q_{d,c}$ is real-rooted, then
\[
 c^2\le\frac8{9(d-2)}.
\]
For its roots $r_i$, write $s_k=\sum r_i^k$. Newton's identities
using the first four coefficients give
\[
 s_1=0,\qquad s_2=2d(d-1),\qquad s_3=-3c d(d-1)(d-2),
 \qquad s_4=4d(d-1)(2d-3).
\]
For real $r_i$ Cauchy--Schwarz on $r_i$ and $r_i^2-s_2/d$ gives
$s_3^2\le s_2(s_4-s_2^2/d)$, which yields the bound. Consequently,
for any $d>2+8/(9c^2)$ and any fixed $n$, $J_{d,n}$ has a nonreal pair
for all sufficiently small $t>0$. Conversely $Q_{d,0}$ is a rescaled
Hermite polynomial with simple real zeros. Thus for any finite set of
degrees and shifts one may first fix a small $c>0$ and then take $t>0$
small enough for every prescribed test to have distinct real roots,
and for one explicitly bounded higher degree to fail. This applies
even with $(\log K_{t,c}(\sqrt{x}))''=-2t^2-6ct^3x<0$ everywhere.
The bound on $d$ is sufficient for detecting failure. Its optimal value
for each $c$ is a further problem.

\paragraph{An exact example.}
At $n=0$, $t=1/10000$, $c=4/5$, the cubic passes and the quartic fails.
The accompanying \texttt{verify\_jensen.py} bounds each $b_k$ by the
degree-nine and degree-ten Taylor polynomials of $e^{-V}$, averaged over
the gamma law with $V=t^2X^2+ct^3X^3$. These are rational bounds because
$\mathbb E X^j=(k+1/2)_j$. In the rescaled variable the cubic has
alternating signs at $-4,0,7/5,3$. For the quartic the determinant
of $(s_{i+j})_{0\le i,j\le2}$ is strictly between $-11000$ and $-1000$.
Real roots would make this matrix a Gram matrix, contradicting that sign.
This finite example and the limiting argument use no zero-height verification.
They describe the reach of a proposed analytic transfer. The zero
distribution of the conductor-$23$ function still requires its own
arithmetic or analytic input.
